\chapter{Revue des outils retenus}
\label{ann:outils}

Les trois critères de sélection sont exposés au chapitre~\ref{ch:entrepot}. Cette annexe en
donne l'application, outil par outil, avec ce qui a été écarté dans chaque cas.

\section{Ingestion : dlt}

Interroger une interface web publique paraît trivial jusqu'à ce qu'on le fasse sérieusement. Il
faut paginer, réessayer sur erreur serveur avec une temporisation croissante, découvrir le schéma
des colonnes, créer les tables, insérer par lots et suivre l'état des chargements.
\textbf{dlt}~\cite{dlt} prend en charge la seconde moitié de cette liste : on déclare une
ressource, la bibliothèque crée et fait évoluer les tables, charge par lots et tient l'historique
des chargements.

La première moitié a dû être écrite. Les particularités des interfaces Hub'Eau décrites au
chapitre~\ref{ch:verrou}, pagination par en-tête de lien, plafonds non documentés, conventions
qui changent d'un point d'accès à l'autre, ne se laissent pas absorber par un mécanisme
générique. Un client dédié d'environ quatre cents lignes porte donc la pagination, les tentatives
successives avec temporisation exponentielle et bruitage, et les garde-fous contre les boucles
infinies. C'est un enseignement en soi : l'outil a économisé le travail répétitif de chargement,
pas l'adaptation aux caprices de la source.

Deux familles d'alternatives ont été écartées. Tout écrire, chargement et gestion des schémas
compris, aurait doublé ce volume sans intérêt propre. Déployer une plateforme d'intégration
complète, du type de celles qui s'administrent par interface graphique, aurait imposé un service
supplémentaire à maintenir et une configuration hors du dépôt, donc non versionnée, pour quatre
sources seulement. dlt occupe l'espace intermédiaire : une bibliothèque, dans le code, versionnée
avec lui.

\section{Transformation : dbt}

Le choix déterminant est ici de transformer \textbf{dans la base} plutôt qu'en mémoire. Avec des
centaines de millions de lignes, charger les données dans un tableau Python n'est pas une
option, alors que le moteur de base de données sait le faire sur place.

\textbf{dbt}~\cite{dbt} est le standard de fait de cette approche, au point d'avoir donné son nom
à un métier, l'\emph{analytics engineering}. Sa maturité et sa diffusion sont précisément ce
qu'on recherche pour un socle destiné à durer. Il apporte quatre mécanismes : chaque
transformation est un fichier SQL versionné, l'ordre d'exécution se déduit automatiquement des
références entre modèles, les tests de qualité se déclarent à côté du modèle qu'ils vérifient,
et la documentation se génère. Un bénéfice secondaire compte autant que les autres : les
transformations restent lisibles par un hydrogéologue ou un analyste, ce qu'un script Python de
plusieurs centaines de lignes n'aurait pas été.

\section{Orchestration : Dagster}

Un ordonnanceur classique raisonne en tâches : exécuter A, puis B, à trois heures du matin.
\textbf{Dagster}~\cite{dagster} raisonne en \emph{actifs} : on déclare qu'une table doit exister
et être à jour, et l'outil détermine ce qu'il faut exécuter pour cela.

La différence est concrète sur cette chaîne. Elle permet de partitionner l'ingestion par année
et de rejouer une seule année, de déclencher la transformation quand la donnée arrive plutôt
qu'à heure fixe, et de voir d'un coup d'œil quelles tables sont fraîches et lesquelles ne le
sont pas. Airflow, l'ordonnanceur historique du domaine, aurait imposé de décrire ces
dépendances à la main et de maintenir un second graphe en parallèle de celui de dbt, alors que
l'intégration est ici native. Un simple ordonnanceur système, quant à lui, n'offre ni reprise
partielle, ni traçabilité, ni limitation de concurrence, trois besoins dont
la suite de ce chapitre montre qu'ils se sont révélés indispensables.

\section{Stockage : PostgreSQL, TimescaleDB et PostGIS}

Le besoin dictait le choix. Ce travail exige simultanément deux capacités : interroger
efficacement des séries temporelles longues, et effectuer des jointures spatiales pour rattacher
chaque station à sa maille climatique. \textbf{PostgreSQL}~\cite{postgresql} est le seul moteur
courant qui réunisse les deux dans un même produit, grâce à \textbf{TimescaleDB}~\cite{timescaledb}
pour le partitionnement temporel et la compression, et à \textbf{PostGIS}~\cite{postgis} pour les
index et opérateurs géospatiaux. PostGIS est d'ailleurs la référence du géospatial libre, employée
bien au-delà du monde de la recherche.

S'y ajoutent trois propriétés qui pesaient lourd dans un contexte de recherche publique : aucun
coût de licence, un déploiement sur le serveur du laboratoire donc aucune dépendance à un
fournisseur ni facturation à la requête, et une interface SQL que tout le monde sait déjà
interroger. Un entrepôt hébergé dans le nuage aurait apporté une élasticité inutile à cette
échelle, en échange d'un budget récurrent et d'une adhérence contractuelle. Un stockage en
fichiers colonnaires, plus léger, ne fournissait ni les jointures spatiales indexées ni la mise
à jour incrémentale dont la chaîne a besoin.

\section{Synthèse des choix}

\begin{table}[H]
\centering
\caption{Synthèse des outils retenus et des alternatives écartées.}
\label{tab:synthese-outils}
\begin{tabularx}{\textwidth}{@{}llX@{}}
\toprule
\textbf{Rôle} & \textbf{Retenu} & \textbf{Écarté, et pourquoi} \\
\midrule
Ingestion & dlt & Tout écrire, chargement et schémas compris ; plateforme d'intégration
administrée (service en plus, configuration hors dépôt) \\
Transformation & dbt & Scripts Python (ne passent pas à l'échelle, illisibles pour les métiers) \\
Orchestration & Dagster & Airflow (dépendances à décrire à la main, second graphe à maintenir) ;
ordonnanceur système (ni reprise, ni traçabilité) \\
Stockage & PostgreSQL & Entrepôt hébergé (coût récurrent, adhérence) ; fichiers colonnaires
(pas de jointure spatiale indexée) \\
Séries temporelles & TimescaleDB & Partitionnement manuel (à écrire et à maintenir) \\
Géospatial & PostGIS & Calculs de distance en Python (lents, non indexés) \\
Déploiement & Docker Compose & Kubernetes (surdimensionné pour un serveur) \\
\bottomrule
\end{tabularx}
\end{table}

Le même raisonnement a présidé aux choix de la plateforme, présentés aux
chapitres~\ref{ch:metier} et~\ref{ch:apprentissage} : Pastas est la référence de la communauté
hydrogéologique pour les modèles à fonction de transfert, Darts unifie sous une seule interface
des architectures de prévision qu'il aurait fallu intégrer une à une, MLflow et Optuna sont les
outils les plus répandus pour, respectivement, tracer des expériences et optimiser des
hyperparamètres. Dans chaque cas, le critère est le même : un successeur doit pouvoir reprendre
le travail sans apprendre un outil que personne n'utilise.

